\chapter{2002:John B. Fenn and Koichi Tanaka and Kurt Wüthrich}

\label{chap:chemistry2002}

The Nobel Prize in Chemistry for the year 2002 was awarded jointly to John B. Fenn and Koichi Tanaka and Kurt Wüthrich.

\textbf{Motivation:}

for their development of soft desorption ionisation methods for mass spectrometric analyses of biological macromolecules (one half, jointly, John B. Fenn and Koichi Tanaka) and for his development of nuclear magnetic resonance spectroscopy for determining the three-dimensional structure of biological macromolecules in solution (the other half, Kurt Wüthrich)

\section{John B. Fenn}

\subsection*{Early Life and Education}

John B. Fenn was born in 1917 in New York City, New York, USA.

\subsection*{Career and Major Research}

Fenn studied at Berea College (B.S., 1937) and earned his Ph.D. from Princeton University in 1940. He was a professor at Yale University from 1967 to 1994, and he later held a position at Virginia Commonwealth University.

\subsection*{Nobel Prize-winning Work}

The work for which John B. Fenn was awarded the Nobel Prize in Chemistry in 2002 was described by the Nobel Committee as: "for their development of soft desorption ionisation methods for mass spectrometric analyses of biological macromolecules".

He developed electrospray ionization, which produces ions from large biomolecules in solution for analysis by mass spectrometry.

\subsection*{Significance and Impact}

Together with the work of Koichi Tanaka, electrospray ionization made mass spectrometry applicable to proteins and other macromolecules.

\subsection*{Later Career and Honors}

Fenn shared the 2002 Nobel Prize in Chemistry, and he died in 2010.

\subsection*{Legacy}

Electrospray ionization is now a standard tool for analyzing proteins and other biomolecules.

\subsection*{Selected Publications}

"Electrospray Ionization for Mass Spectrometry of Large Biomolecules" (Science, 1989)

\clearpage

\section{Koichi Tanaka}

\subsection*{Early Life and Education}

Koichi Tanaka was born in 1959 in Toyama, Japan.

\subsection*{Career and Major Research}

Tanaka joined the Shimadzu Corporation in 1983 as a research engineer, and he remained with the company.

\subsection*{Nobel Prize-winning Work}

The work for which Koichi Tanaka was awarded the Nobel Prize in Chemistry in 2002 was described by the Nobel Committee as: "for their development of soft desorption ionisation methods for mass spectrometric analyses of biological macromolecules".

He developed soft laser desorption mass spectrometry, using a laser with a metal powder and glycerol matrix to ionize large intact biomolecules.

\subsection*{Significance and Impact}

His method allowed proteins and other macromolecules to be analyzed without fragmentation, making mass spectrometry useful for biology.

\subsection*{Later Career and Honors}

Tanaka remained at the Shimadzu Corporation. He shared the 2002 Nobel Prize in Chemistry.

\subsection*{Legacy}

His approach contributed to the methods that now underlie the mass spectrometry of biological macromolecules.

\subsection*{Selected Publications}

"Protein and Polymer Analyses up to m/z 100 000 by Laser Ionization Time-of-Flight Mass Spectrometry" (Rapid Communications in Mass Spectrometry, 1988)

\clearpage

\section{Kurt Wüthrich}

\subsection*{Early Life and Education}

Kurt Wüthrich was born in 1938 in Aarberg, Switzerland.

\subsection*{Career and Major Research}

Wüthrich studied at the University of Bern and earned his Ph.D. from the University of Basel in 1964. He was a professor at the ETH Zurich from 1969, and he also held a position at the Scripps Research Institute in La Jolla.

\subsection*{Nobel Prize-winning Work}

The work for which Kurt Wüthrich was awarded the Nobel Prize in Chemistry in 2002 was recognized for his development of nuclear magnetic resonance spectroscopy for determining the three-dimensional structure of biological macromolecules in solution.

He developed methods for determining complete three-dimensional structures of proteins in solution by NMR, using the nuclear Overhauser effect to measure distances between atoms.

\subsection*{Significance and Impact}

The methods made NMR a second technique, alongside X-ray crystallography, for protein structure determination.

\subsection*{Later Career and Honors}

Wüthrich continued his research at the ETH Zurich and the Scripps Research Institute. He shared the 2002 Nobel Prize in Chemistry.

\subsection*{Legacy}

His structural NMR methods are now used routinely to study proteins in solution.

\subsection*{Selected Publications}

"NMR of Proteins and Nucleic Acids" (Wiley, 1986)

\clearpage

\section*{Chapter Summary}

The 2002 Nobel Prize in Chemistry recognized methods for analyzing biological macromolecules. Fenn and Tanaka developed ionization methods for mass spectrometry, while Wüthrich developed NMR for structure determination in solution.
